\section{Introduction}

Ce document recense et décrit les \textbf{données échangées} entre l'automate
Schneider M340 et les différents organes de la cellule robotique : le poste
guichet du convoyeur (via le contrôleur WAGO PFC200), la baie de commande du
robot SCARA (contrôleur Stäubli CS9, serveur \texttt{uniVALplc}) et le pupitre
opérateur / écran de supervision (contrôleur Schneider M221). Pour la
présentation générale du système et de ses contraintes, se reporter au dossier
\texttt{01\_conception} ; pour la description fonctionnelle des modes de
marche et d'arrêt, se reporter au dossier \texttt{02\_analyse\_fonctionnelle}.

\subsection{Objet du document}

Contrairement à l'analyse fonctionnelle des modes, ce document a une vocation
délibérément \textbf{technique} : il décrit précisément la nature, le format
et le contenu des données échangées avec chacun des équipements de la
cellule, telles qu'elles sont définies par les fichiers de configuration
EDS (Electronic Data Sheet) et les structures de données fournies par les
bibliothèques logicielles utilisées (\texttt{uniVALplc} côté robot,
assemblage Ethernet/IP côté PFC200 et M221).

\subsection{Rappel du réseau Ethernet/IP}

L'automate M340 est équipé de deux coupleurs Ethernet/IP de type
BMX~NOC~0401.2, chacun scrutant un ensemble d'équipements distinct :

\begin{itemize}
    \item \textbf{NOC\_ETHIP\_CONVOYEUR} : scrute le contrôleur
    \textbf{PFC200} (WAGO), qui gère la commande du poste guichet du
    convoyeur ;
    \item \textbf{NOC\_ETHIP\_ROBOT} : scrute conjointement la baie
    \textbf{CS9} du robot SCARA (serveur \texttt{uniVALplc}) et le
    contrôleur \textbf{M221} (îlot d'entrées/sorties du pupitre opérateur et
    interface de supervision Modbus-TCP vers l'écran HMI).
\end{itemize}

\begin{UPSTIinfor}{Vocabulaire Ethernet/IP}
Pour chaque équipement scruté par l'automate (le \emph{Scanner} /
\emph{Originator}), on distingue :
\begin{itemize}
    \item les \textbf{données produites} (\texttt{T->O}, \emph{Target ->
    Originator}) : elles vont de l'équipement (la \emph{Target} /
    \emph{Adapter}) vers l'automate. Elles constituent, du point de vue de
    l'automate, une \textbf{entrée} (\texttt{IN}) ;
    \item les \textbf{données consommées} (\texttt{O->T}, \emph{Originator
    -> Target}) : elles vont de l'automate vers l'équipement. Elles
    constituent, du point de vue de l'automate, une \textbf{sortie}
    (\texttt{OUT}).
\end{itemize}
\end{UPSTIinfor}

Chaque coupleur réserve, dans l'espace mémoire \texttt{\%MW} de l'automate,
une zone d'entrée (contenant l'état de santé des connexions --
\texttt{HEALTH\_BITS\_IN} -- suivi de la zone image des données produites)
et une zone de sortie (contenant la mise en service des connexions --
\texttt{CONTROL\_BITS\_OUT} -- suivie de la zone image des données
consommées).

Les sections suivantes détaillent, pour chaque organe de la cellule, la
nature des données échangées.
